\documentclass[10pt,a4paper]{amsart}
\usepackage[margin=19mm]{geometry}
\usepackage{amsmath,amssymb,mathtools}
\usepackage[scr=rsfs,cal=euler]{mathalfa}
\usepackage{xfrac}
\usepackage[unicode,hidelinks,bookmarks=false]{hyperref}
\newcommand{\ie}{i.e.\ }
\newcommand{\cf}{cf.\ }
\newcommand{\resp}{respectively\ }
\newcommand{\Subsection}{\subsection}
\providecommand{\nobreakdash}{\nobreak\hbox{-}}
\setlength{\emergencystretch}{2em}
\multlinegap0pt
\usepackage{amssymb}
\usepackage{mathtools}
\usepackage{tikz}
\usepackage{mdframed}
\usepackage{tcolorbox}
\usepackage{enumerate}
\usepackage[shortlabels]{enumitem}
\newtheorem{lemma}{Lemma}
\usepackage{cleveref}
\crefname{lemma}{Lemma}{Lemmas}
\newcommand{\eps}{\varepsilon}
\newcommand{\mb}{\mathbb}
\newcommand{\on}{\operatorname}
\title{Permanents of random matrices over finite fields}
\author{Zach Hunter, Matthew Kwan, Lisa Sauermann}
\date{Source: arXiv:2603.15856v1 (CC BY 4.0)}
\begin{document}
\maketitle

\noindent\emph{Source and licence.} Permanents of random matrices over finite fields. Version arXiv:2603.15856v1 (CC BY 4.0), distributed by arXiv under the Creative Commons Attribution 4.0 International licence, http://creativecommons.org/licenses/by/4.0/, as recorded in the arXiv licence metadata for this exact version and shown on the abstract page https://arxiv.org/abs/2603.15856v1. Full licence notice: \texttt{licenses/arxiv-2603.15856-CC-BY-4.0.txt}.\par\smallskip

\noindent\textit{Editorial scope.} Counted unit: the article's lemma lem:growth (source0.tex lines 575--582). The article's complete original proof of it is delivered with it (source0.tex lines 584--655, one \texttt{proof} environment of 72 lines). No mathematics is written, repaired or re-derived: the statement and the proof are the article's own lines, in the article's own notation, and no retained line is substituted or re-flowed.  Retained uncounted as the source-local context the proof needs: lines 307--313.  Nothing was omitted from any retained span, so the package carries no omission record.  No retained line contains a citation, so the wrapper carries no bibliography and no bibliography entry.  No retained line contains a figure environment, an image-inclusion command or drawing code, so no image is delivered and the package declares no asset dependency.  Editorial formatting only, in addition: an AMS \texttt{amsart} preamble that restates the article's own declarations and macros used by the retained lines, namely the environments \texttt{lemma} on separate counters, together with the standard packages those lines need.  The retained environments print in this document as Lemma 1 in order of appearance, whereas the article prints the same items with the numbering of its own full document.  The complete native source of the article as distributed in the arXiv e-print for this exact version is bundled unchanged as \texttt{sources/196504/source0.tex}.
\begin{lemma}
\label{lem:match-det-0}For any $s\in\{1,\dots,n\}$, and any outcome
of $A^{\uparrow s}$ satisfying $\mathcal{E}(s)$, we have
\[
\Pr[\mathcal{E}(s-1)\,|\,A^{\uparrow s}]\ge1-q^{-s}.
\]
\end{lemma}

\begin{lemma}\label{lem:growth}
For any $\varepsilon>0$, there is $T=T_\mu(\varepsilon)\in\mb N$ such that the following
holds. Consider any subset $J\subseteq\{1,\dots,n\}$ of size $|J|=2T$.
Then,
\[
\Pr[\text{there is a}\text{ subset }I\subseteq J\text{ of size }|I|=T\text{ with }\on{per}(A;I)\ne0]\ge 1-\varepsilon.
\]
\end{lemma}

\begin{proof} Let $\rho=\max_{z\in\mb F_{p}}\mu(z)<1$ (recalling that $\mu$
is supported on at least two values), and choose $0<\delta<1$
small enough such that $1-\rho(1+\delta)>\delta$. Furthermore choose $T>2(1+\delta)/(\eps\delta^2)$ to be sufficiently large such that $\rho^{\delta T}/(1-\rho)<\eps/4$.

We consider a random process that attempts to construct a nested sequence
of sets $I_{0}\supseteq\dots\supseteq I_{n-T}$ (starting with $I_{0}=\{1,\dots,n\}$,
and removing an element at each step), in such a way that $\on{per}(A;I_{t})\ne0$
for all $t=0,\dots,n-T$, and $I_{n-T}\subseteq J$. If this process succeeds,
then we can take $I_{n-T}$ as the desired set $I$.

Specifically, the process is defined as follows. Let $T'=2T+\lfloor \delta T\rfloor$. For each step $t=1,\dots,n-T'$, we define $I_t$ as follows:
\begin{itemize}
\item If there is $i\in I_{t-1}\setminus J$ such that $\on{per}(A;I_{t-1}\setminus\{i\})\ne0$, take the minimum such $i$ and
set $I_{t}=I_{t-1}\setminus\{i\}$.
\item Otherwise, terminate the process.
\end{itemize}
For each step $t=n-T'+1,\dots, n-T$, we define $I_t$ as follows:
\begin{itemize}
\item If there is $i\in I_{t-1}\setminus J$ such that $\on{per}(A;I_{t-1}\setminus\{i\})\ne0$, take the minimum such $i$ and
set $I_{t}=I_{t-1}\setminus\{i\}$.
\item Otherwise, if there is any $i\in I_{t-1}$ such that $\on{per}(A;I_{t-1}\setminus\{i\})\ne0$, take the minimum such $i$ and
set $I_{t}=I_{t-1}\setminus\{i\}$.
\item Otherwise, terminate the process.
\end{itemize}
That is to say: in each of the first $n-T'$ steps we make sure to
remove an element outside $J$. In each of the subsequent
steps we still try to remove an element outside $J$ if possible,
but if not we allow ourselves to remove some other element. Note that we always have $|I_t|=n-t$ and $\on{per}(A;I_{t})\ne0$ whenever the set $I_t$ is defined, and also note that $I_t$ only depends on the outcome of $A^{\uparrow (n-t)}$ (i.e., only on the first $t$ rows of the matrix $A$).

It suffices to show that with probability at least $1-\varepsilon/2$
this process does not terminate in any step before defining $I_{n-T}$, and to show
that\footnote{If we terminate in some step (so $I_{n-T}$ is not defined), we say that the event $\{I_{n-T}\not \subseteq J\}$ does \emph{not} occur.} $I_{n-T}\nsubseteq J$ with probability at most $\varepsilon/2$.

First, note that for any outcome
of $A^{\uparrow (n-t)}$ for which $\on{per}(A;I_{t})\ne0$, we have
\[
\Pr[\text{the process terminates at step }t+1\,|\,A^{\uparrow (n-t)}]\le\begin{cases}
\rho^{|I_{t}\setminus J|}\le\rho^{n-t-2T} & \text{if }t<n-T'\\
\rho^{|I_{t}|}=\rho^{n-t}& \text{if }n-T'\le t< n-T
\end{cases}
\]
The reasoning is basically the same as in the proof of \cref{lem:match-det-0}. Indeed,
write $A=(a_{i,j})$, condition on an outcome of $A^{\uparrow (n-t)}$,
and additionally condition on $a_{t+1,i}$ for all $i\notin I_{t}$.
Then, the quantities $\on{per}(A;I_{t}\setminus\{i\})$, for
$i\in I_{t}$, become independent random variables (each only depending on the outcome of $a_{t+1,i}$), which each have
probability at most $\rho$ of taking any particular value (in particular,
of taking the value zero).

By a union bound, it follows that the probability that the process terminates in some step (i.e., the probability that $I_{n-T}$ does not get defined) is at most

\[
\sum_{t=0}^{n-T'-1}\rho^{n-t-2T}+\sum_{t=n-T'}^{n-T-1}\rho^{n-t}=\sum_{j=T'+1-2T}^{n-2T}\rho^{j}+\sum_{j=T+1}^{T'}\rho^{j}\le \frac{\rho^{\lfloor \delta T\rfloor+1}+\rho^{T+1}}{1-\rho}\le \frac{2\rho^{\delta T}}{1-\rho}\le\frac{\varepsilon}{2}.
\]

Now we study the probability that $I_{n-T}\nsubseteq J$. For $t= n-T'+1,\dots,n-T$, say step $t$ is \emph{bad} if $I_{t}\setminus J=I_{t-1}\setminus J\ne \emptyset$ (i.e., if $I_{t-1}\nsubseteq J$, and the element we remove from $I_{t-1}$ in step $t$ lies in $J$). Then for any outcome of $A^{\uparrow t}$ with $I_{t-1}\nsubseteq J$, we have
\[
\Pr[t\text{ is bad}\,|\,A^{\uparrow t}]\le\rho^{|I_{n-t}\setminus J|}\le\rho,
\]
and for any outcome of $A^{\uparrow t}$ with $I_{t-1}\subseteq J$ we have $\Pr[t\text{ is bad}\,|\,A^{\uparrow t}]=0\le \rho$. Thus, the number of bad steps is stochastically dominated by a $\on{Binomial}(T'-T,\rho)$-distributed random variable.

Now, if $I_{n-T}\not\subseteq J$, there must be an element $h\in I_{n-T}\not\subseteq J$. This means that $I_{t}\not\subseteq J$ for all $t= n-T'+1,\dots,n-T$. Furthermore, at most $|I_{n-T'}\setminus J|\le |I_{n-T'}|-|J|=T'-2T=\lfloor \delta T\rfloor$ elements outside $J$ get removed during the $T'-T$ steps $n-T'+1,\dots,n-T$. This means that at least $T'-T-\lfloor \delta T\rfloor=T$ of these steps must be bad.

Thus, writing $B\sim\on{Binomial}(T'-T,\rho)$ (which has expectation $\rho(T'-T)$ and standard deviation at
most $\sqrt{T'-T}\le \sqrt{(1+\delta)T}$), we have
\begin{align*}
\Pr[I_{n-T}\not\subseteq J] & \le\Pr[\text{at least }T\text{ of the steps }n-T'+1,\dots,n-T\text{ are bad}]\\
 & \le\Pr[B\ge T]\\
 & \le\Pr[B\ge\rho(T'-T)+\delta T]\le\frac{(1+\delta)T}{\delta^2 T^2}=\frac{(1+\delta)}{\delta^2 T}\le \frac{\varepsilon}{2},
\end{align*}
noting that $\rho(T'-T)+\delta T\le \rho(1+\delta) T+\delta T\le T$ by definition of $T'$ and of $\delta$, and using Chebyshev's inequality.
\end{proof}

\end{document}
